% Question 2 -- Motion or Memory?  (included by main.tex)
\setcounter{table}{0}
\setcounter{figure}{0}
\newcommand{\qa}[1]{\paragraph{Q2.#1}}
\newcommand{\Ex}[1]{E^{(#1)}}

\paragraph{Conventions.} Every table entry for a learned model is \emph{mean (seed 0, seed 1)};
Blank and Copy-last are deterministic and get a single value. Errors have three significant figures,
and values below $10^{-3}$ use scientific notation. All model results are on the 1,000 test videos
(9,000--9,999) at $t=1$, unless stated otherwise. Code: \texttt{Question2/} in the repository
(\texttt{Question2/README.md} lists the command that produces each table and figure).

\paragraph{Hardware and timing.} Everything was run on one laptop CPU: AMD Ryzen~5 7430U
(6 cores, 12 threads), 15\,GB RAM, no GPU, Linux 7.2, Python~3.14, PyTorch~2.15.0 (CPU build).
Training used 12 threads, with one run at a time. Inference was timed with PyTorch's default 6 threads,
following the protocol in the handout (batch 32, \texttt{eval()}, \texttt{no\_grad()}, all 1,000 test videos).
\emph{Note on model~A, seed~0:} its first five epochs (303, 276, 225, 221 and 227\,s) ran while other
applications and package installation shared the CPU, and training was resumed from checkpoints after
interruptions (resuming reproduces the uninterrupted weights bit for bit). Its reported 30-epoch mean of
210\,s is the measured value. Its epochs 6--30 average 202\,s, the same as the other baseline-architecture
runs (202--205\,s). The same applies to \emph{model~C, seed~0}: its epochs 4--13 (220--275\,s, mean 242\,s)
shared the CPU with other work and two interruptions, so its measured mean is 216\,s; its other 20 epochs
average 203\,s. Model~D and D+B take about 60\,s per epoch.

%======================================================================
\section{Setup}
%======================================================================
The videos are split by dataset index (0--7,999 train, 8,000--8,999 validation, 9,000--9,999 test)
and scaled to $[0,1]$. All checks in Table~\ref{tab:q2-t1} call the same transformation and window
functions that the experiments use (\texttt{q2/data.py}), and every row matches.

\begin{table}[H]
\centering\small
\caption{Checks of the transformation and window functions.}
\label{tab:q2-t1}
\begin{tabular}{lll}
\toprule
Check & Expected & Yours \\
\midrule
(a) max$|$M(M(x)) - x$|$ & 0 & 0 (torch.equal=True) \\
(b) max$|$R$^4$(x) - x$|$ & 0 & 0 (torch.equal=True) \\
(c) max$|$M R M(x) - R$^{-1}$(x)$|$ & 0 & 0 (torch.equal=True) \\
(d) bright pixel of M(p) & (7, 10, 43) & (7, 10, 43) \\
(e) bright pixel of R(p) & (7, 43, 10) & (7, 43, 10) \\
(f) 1x input / target frames & 1-5 / 6-10 & 1, 2, 3, 4, 5 / 6, 7, 8, 9, 10 \\
(g) 2x input / target frames & 1, 3, 5, 7, 9 / 11, 13, 15, 17, 19 & 1, 3, 5, 7, 9 / 11, 13, 15, 17, 19 \\
(h) 3x input / target frames & 1, 4, 7, 10, 13 / 16, 19 & 1, 4, 7, 10, 13 / 16, 19 \\
\bottomrule
\end{tabular}

\end{table}

\begin{figure}[H]
\centering
\includegraphics[width=\linewidth]{fig1_flow.png}
\caption{Farneback optical flow on the first 50 training videos (950 consecutive frame pairs).
103,050 pixels are kept (intensity $>0.2$ in both frames, flow magnitude $>0.5$\,px).
(a) Flow direction $\theta=\mathrm{atan2}(-d_y,d_x)$ in 18 bins of $20^\circ$, as the fraction of kept pixels;
the dashed line is the uniform level $1/18$.
(b) Flow speed in px/frame, as the fraction of kept pixels; the median is 3.32\,px/frame (dashed line).}
\label{fig:q2-f1}
\end{figure}

\qa{1} \emph{Directions.} Every $20^\circ$ bin holds between 2.6\,\% and 8.2\,\% of the kept pixels
(uniform: 5.6\,\%). No direction is missing; the unevenness comes from having only about 100 digit
trajectories (50 videos, two digits each). A mirror maps $\theta\mapsto180^\circ-\theta$, and a $90^\circ$
rotation maps $\theta\mapsto\theta+90^\circ$. So mirroring or rotating every training video would only
reorder the bars of Figure~\ref{fig:q2-f1}(a). Applying these maps to the measured histogram moves it by a
total-variation distance of 0.132 (mirror) and 0.200 (rotation), comparable to its distance from uniform
(0.110). Every direction in a mirrored or rotated video is therefore a direction the model has seen in
at least 2.6\,\% of its training pixels, so the trajectories themselves are familiar. Only their
frequencies change, and rotation reshuffles them more than mirroring.
\emph{Speeds.} The median speed is 3.32\,px/frame (interquartile range 3.00--3.89, 95th percentile 4.32,
99th 4.67). At 2$\times$ the digits move about 6.6\,px per observed frame (95th percentile 8.6), and at
3$\times$ about 10.0\,px (95th percentile 13.0). Both are far outside the training range, where almost no
pixel moves more than 5\,px.
\emph{Overlap.} Farneback assumes that a pixel keeps its brightness as it moves and that motion is locally
smooth. Where two digits overlap, both assumptions fail: the pixel saturates, and two motions share one
neighbourhood. The flow there is a blend of the two directions and is usually too short, so it adds a
little mass to the slow tail of Figure~\ref{fig:q2-f1}(b) (with bounces and digit edges) and smears a few
pixels across directions. Overlaps cover a small fraction of the kept pixels and do not remove any
direction, so they can slightly lower the median speed and smooth the histogram, but cannot change
either answer above.

\qa{2} \emph{Mirror}: changes the appearance (a mirrored 2 is not a 2) and relabels directions
($\theta\mapsto180^\circ-\theta$), with speeds unchanged; every resulting direction occurs in training.
\emph{Rotation}: changes the appearance even more (digits lie on their side, so vertical and horizontal
strokes swap) and relabels directions ($\theta\mapsto\theta+90^\circ$), again with speeds unchanged.
\emph{Speed-up}: leaves the appearance untouched and changes only the motion statistics. Every
displacement is multiplied by $k$, beyond anything seen in training, and bounces fall between observed
frames more often. Together they form a controlled experiment. The true dynamics are invariant to all
three, and the correct transformed future $T(Y)$ is known exactly, so any failure belongs to the model,
not the task. Each transformation also changes a different factor: speed-up isolates unfamiliar
\emph{motion} at fixed appearance, while mirror and rotation keep the motion within the training
distribution and mainly change \emph{appearance}, by a smaller (mirror) and larger (rotation) amount.
Comparing the three therefore indicates which factor a failure depends on.

%======================================================================
\section{Baseline: ConvLSTM}
%======================================================================
Model A is implemented exactly as specified (\texttt{q2/models.py}): two peephole-free ConvLSTM
layers with 64 hidden channels and $3\times3$ kernels on $4\times4$-folded frames, a $1\times1$ output
convolution, \texttt{pixel\_shuffle} and a sigmoid, and a 9-step closed-loop rollout without teacher
forcing. It has exactly 480,784 trainable parameters. Training follows the protocol in the handout
(\texttt{q2/train.py}).

\begin{table}[H]
\centering
\caption{Per-horizon test error, $\Ex5$ and parameter count of model A and the two reference predictors (1$\times$, $t=1$).}
\label{tab:q2-t2}
\resizebox{\linewidth}{!}{\begin{tabular}{lccccccc}
\toprule
Model & MSE$_1$ & MSE$_2$ & MSE$_3$ & MSE$_4$ & MSE$_5$ & $E^{(5)}$ & Params \\
\midrule
A & 0.0172 (0.0175, 0.0169) & 0.0206 (0.0210, 0.0202) & 0.0230 (0.0234, 0.0225) & 0.0247 (0.0253, 0.0242) & 0.0264 (0.0270, 0.0259) & 0.0224 (0.0228, 0.0219) & 480,784 \\
Blank & 0.0426 & 0.0426 & 0.0426 & 0.0425 & 0.0425 & 0.0426 & 0 \\
Copy-last & 0.0492 & 0.0604 & 0.0669 & 0.0708 & 0.0732 & 0.0641 & 0 \\
\bottomrule
\end{tabular}
}
\end{table}

\begin{figure}[H]
\centering
\includegraphics[width=0.62\linewidth]{fig2_model_A_examples.png}
\caption{Model A (seed 0) on test videos 9007, 9012 and 9851: input frames 1--5, predictions and
ground truth for frames 6--10. In 9851 the 3 and the 1 touch in frames 4--5, overlap in frames 6--8
and then separate.}
\label{fig:q2-f2}
\end{figure}

\qa{3} \emph{Reach of model A.} Folding $4\times4$ patches makes one cell of the $16\times16$ grid
equal to 4 original pixels. Only the two $3\times3$ convolutions move information, each by at most one
cell (4\,px) along a row or column per application. Within step $\tau$ there are three kinds of path to
$\hat x_{\tau+1}$:
(i) from the new input $u_\tau$: layer~1's convolution (one cell) gives $h^{(1)}_\tau$, and layer~2's
convolution (one more cell) gives $h^{(2)}_\tau$, which the $1\times1$ convolution and \texttt{pixel\_shuffle}
turn into the frame without moving anything: 2 cells $=8$\,px;
(ii) from $h^{(1)}_{\tau-1}$: it enters layer~1's convolution together with $u_\tau$, so it also crosses
both convolutions: 8\,px;
(iii) from $h^{(2)}_{\tau-1}$: it enters only layer~2's convolution: 4\,px.
The longest path gives a reach of \textbf{8\,px per time step} along a row or column. Along a diagonal,
the $3\times3$ neighbourhood moves one cell in both directions, so it is $8\sqrt2\approx11.3$\,px.
The recurrence does not lengthen the path within one step: it only carries what earlier steps already
moved. Over several steps, though, the reach adds up.

\qa{4} Model A beats both reference predictors at every horizon. Its $\Ex5=0.0224$ is 47\,\% below
Blank (0.0426) and 65\,\% below Copy-last (0.0641), and the two seeds differ by only 4\,\%. Copy-last is
worse than Blank: a digit drawn in the wrong place is penalised twice (once where it is, once where it
should be), whereas a blank frame is penalised only once. Model A's error grows steadily, from
$\mathrm{MSE}_1=0.0172$ to $\mathrm{MSE}_5=0.0264$ ($1.5\times$), while Blank stays flat and Copy-last
rises from 0.0492 to 0.0732 as the digits move away from their last position. The growth comes from
compounding blur. Each prediction is fed back as the next input, so uncertainty about position and
shape builds up, and the MSE-trained network hedges by spreading intensity. In Figure~\ref{fig:q2-f2}
the digits are recognisable at $h=1$ but have become grey, textured blobs by $h=4$--5, which sit roughly
where the digits are.
\emph{Overlap.} In 9851, the 3 and the 1 merge in the last input frames. The prediction keeps a single
merged blob that stays near the point of contact, spreads into a U-shaped smear and fades. It never
separates the two digits again, although in the truth they pass through each other and re-emerge as
``13''. Once the digits overlap, the model has no separate memory of each digit's identity and velocity,
so it predicts an average.

\qa{5} \emph{The answer below is included exactly as first written. It was committed to the repository
(\texttt{Question2/predictions/Q2\_5.md}, commit \texttt{bfc21b6}, 2026-10-27 17:02 PKT) before any
experiment of Section~3 was run (the first one ran on 2026-10-29 at 04:24 PKT).}

\begin{quote}\small
\textbf{Q2.5: Predictions before running ``Testing the Baseline''}

Written before any experiment of Section 3 (mirror, rotation, speed-up) was run on model A.
Inputs to this prediction: the reach estimate of Q2.3 and Figure 1 only.

\textbf{Reach of model A (Q2.3).} One 3x3 convolution on the 16x16 folded grid moves information one
cell = 4 original pixels. Within one time step the longest path is u\_tau (or h\^{}(1)\_\{tau-1\}) $\to$
conv of layer 1 $\to$ h\^{}(1)\_tau $\to$ conv of layer 2 $\to$ h\^{}(2)\_tau $\to$ 1x1 output, i.e.\ two 3x3 hops = 8 px
along a row or column (8*sqrt(2) $\sim$ 11.3 px along a diagonal). The path from h\^{}(2)\_\{tau-1\} only
crosses one 3x3 convolution (4 px).

\textbf{Figure 1.} Every 20-degree direction bin holds 2.6-8.2 \% of the moving pixels (uniform = 5.6 \%),
so all directions occur in training although not equally often. The median speed is 3.32 px/frame
(95th percentile 4.3 px/frame); 2x is therefore about 6.6 px/frame and 3x about 10 px/frame.

\textbf{Predictions.}
\begin{enumerate}[nosep]
\item \emph{Mirror test:} model A will largely pass. Mirrored trajectories are motions it has seen (the
   direction histogram is not uniform but covers every direction), and constant-velocity motion
   with bounces is a local rule that a convolutional recurrence can apply anywhere. The mirrored
   digits are unfamiliar shapes, so I expect a small but clearly nonzero consistency gap C\_M (a
   fraction of E(5), not zero) and E\_M only slightly above E(5): blurrier, less recognisable
   digits, but moving in the right direction.
\item \emph{Rotation test:} also mostly passes, but somewhat worse than mirror: rotation changes the
   digit appearance more (digits lie on their side) and maps the direction histogram onto a
   less similar one, so C\_R $>$ C\_M and E\_R $>$ E\_M.
\item \emph{Speed:} at 2x (about 6.6 px/frame, up to about 8.6 px at the 95th percentile) the motion is
   mostly still within the 8 px reach, so model A degrades (these speeds were never seen in
   training) but still tracks the digits. At 3x (about 10 px/frame, well above 8 px along an
   axis) the digit moves further than information can travel in one update, so the model should
   break down: digits will stall, blur into a smear or vanish instead of following the motion.
   I expect E(2) at 3x to be close to or worse than Copy-last.
\end{enumerate}
\end{quote}

%======================================================================
\section{Testing the Baseline}
%======================================================================
\begin{table}[H]
\centering
\caption{Accuracy at 1$\times$ and consistency under mirror ($M$) and rotation ($R$).}
\label{tab:q2-t3}
\resizebox{\linewidth}{!}{\begin{tabular}{lccccc}
\toprule
Model & $E^{(5)}$ & $C_M$ & $E_M$ & $C_R$ & $E_R$ \\
\midrule
A & 0.0224 (0.0228, 0.0219) & 0.00156 (0.00148, 0.00164) & 0.0230 (0.0234, 0.0226) & 0.00275 (0.00243, 0.00307) & 0.0240 (0.0243, 0.0238) \\
Blank & 0.0426 & 0 & 0.0426 & 0 & 0.0426 \\
Copy-last & 0.0641 & 0 & 0.0641 & 0 & 0.0641 \\
\bottomrule
\end{tabular}
}
\end{table}

\begin{figure}[H]
\centering
\begin{subfigure}{\linewidth}\centering
\includegraphics[width=0.92\linewidth]{fig3_A_M_9269.png}\caption{$T=M$, video 9269 (largest per-video $C_M$)}
\end{subfigure}
\begin{subfigure}{\linewidth}\centering
\includegraphics[width=0.92\linewidth]{fig3_A_M_9256.png}\caption{$T=M$, video 9256 (second largest $C_M$)}
\end{subfigure}
\caption[Model A, largest consistency gaps]{Model A (seed 0), the two test videos with the largest per-video $C_M$ (continued below for $C_R$).
Rows: (1) $X\,|\,f(X)$; (2) $T(X)\,|\,T(Y)$; (3) $T(X)\,|\,f(T(X))$; (4) $T(X)\,|\,T(f(X))$.
Animated GIFs that step through the five horizons:
\path{Question2/gifs/fig3_modelA_M_video9269.gif}, \path{Question2/gifs/fig3_modelA_M_video9256.gif},
\path{Question2/gifs/fig3_modelA_R_video9983.gif}, \path{Question2/gifs/fig3_modelA_R_video9733.gif}.}
\label{fig:q2-f3}
\end{figure}

\begin{figure}[H]
\ContinuedFloat
\centering
\begin{subfigure}{\linewidth}\centering
\includegraphics[width=0.92\linewidth]{fig3_A_R_9983.png}\caption{$T=R$, video 9983 (largest per-video $C_R$)}
\end{subfigure}
\begin{subfigure}{\linewidth}\centering
\includegraphics[width=0.92\linewidth]{fig3_A_R_9733.png}\caption{$T=R$, video 9733 (second largest $C_R$)}
\end{subfigure}
\caption{(continued) The two test videos with the largest per-video $C_R$.}
\end{figure}

\begin{figure}[H]
\centering
\includegraphics[width=\linewidth]{fig4_consistency_hist_A.png}
\caption{Per-video consistency gaps of model A (seed 0) over the 1,000 test videos: (a) $C_M$,
(b) $C_R$. Red dashed lines mark the Figure~\ref{fig:q2-f3} videos.}
\label{fig:q2-f4}
\end{figure}

\qa{6} Blank and Copy-last commute with every pixel permutation, so $C_M=C_R=0$ exactly. Yet their
transformed errors equal their ordinary errors (0.0426 and 0.0641), far worse than model A's $E_M=0.0230$.
A forecaster that ignores the motion, or ignores its input altogether, is perfectly consistent. So $C_T$
on its own measures only whether a model treats a transformed scene the same way as the original, not
whether that treatment is right. A small gap is necessary for having learned a transferable motion rule,
but it is not sufficient. $E_T$ checks that the consistent prediction is also correct. Reading $C_T$ next
to $E_T$ separates ``consistent and right'' (small $C_T$, small $E_T$) from ``consistent and useless''
(Blank) and from ``inconsistent'' (large $C_T$).

\qa{7} \emph{Video 9983, $T=R$} (Figure~\ref{fig:q2-f3}c). In the original video, a 3 and a 0 overlap
in all five input frames. After the rotation they are stacked vertically, and in the truth (row 2) they
separate as they move: by $h=5$ the rotated ``30'' stands clearly apart near the lower left.
Row~3 ($f(R(X))$) shows what the model does when it sees the rotated input. It keeps the pair as one blob
but stretches it \emph{downwards} into a tall streak, so it follows the true downward drift of the
lower digit, at least partly. Row~4 ($R(f(X))$) is the model's account of the original video, rotated. Here
the merged blob stays compact and square, with a checkerboard texture, close to where the inputs ended.
It drifts slightly but never stretches downwards, and fades in place. So the two accounts disagree on
where the mass goes: one moves part of it along the true direction, the other leaves it where it was.
Neither recovers the two separate digits; in both, the pair blurs into something unrecognisable. The
largest gaps belong to videos like this one, with overlapping, curly digits, where appearance is
ambiguous and the model's motion estimate depends on the shape of the blob it sees.

\qa{8} Yes, rotation is harder for model A. $C_R=0.00275$ is $1.8\times$ $C_M=0.00156$ (in both
seeds), and $E_R-\Ex5=0.0016$ is 2.7 times $E_M-\Ex5=0.0006$; in Figure~\ref{fig:q2-f4} the whole $C_R$
distribution sits to the right of $C_M$. By Q2.1 and Q2.2, both transformations keep every speed and map
every direction onto one that occurs in at least 2.6\,\% of the training pixels; the rotated direction
histogram is only somewhat less similar (total variation 0.200 vs 0.132). Rotation, however, changes
the appearance much more: digits lie on their side, and the vertical-stroke statistics that the $3\times3$
filters on $4\times4$ patches are tuned to become horizontal ones, whereas mirroring keeps the strokes
upright and leaves several digits (0, 1, 8) almost unchanged. The difference is therefore mostly
\emph{appearance}, not motion. This is consistent with where the largest gaps occur (Q2.7): overlapping
or unusually shaped digits, where the model's motion estimate depends on how the blob looks. The direction
reweighting may add a smaller part, which these tests alone cannot fully separate from appearance.

\begin{table}[H]
\centering
\caption{Model A at faster playback ($t=1$), with Copy-last for comparison. $\Ex2$ uses the first two horizons only.}
\label{tab:q2-t4}
\resizebox{\linewidth}{!}{\begin{tabular}{lccccc}
\toprule
Model & $E^{(2)}$ 1$\times$ & $E^{(2)}$ 2$\times$ & $E^{(2)}$ 3$\times$ & $E^{(5)}$ 1$\times$ & $E^{(5)}$ 2$\times$ \\
\midrule
A & 0.0189 (0.0193, 0.0185) & 0.0439 (0.0428, 0.0450) & 0.0513 (0.0505, 0.0522) & 0.0224 (0.0228, 0.0219) & 0.0476 (0.0459, 0.0493) \\
Copy-last & 0.0548 & 0.0661 & 0.0703 & 0.0641 & 0.0719 \\
\bottomrule
\end{tabular}
}
\end{table}

\begin{figure}[H]
\centering
\includegraphics[width=\linewidth]{fig5_speed_A.png}
\caption{Model A (seed 0) on test video 9021 at 1$\times$, 2$\times$ and 3$\times$ ($t=1$): the five input
frames, the first two predictions and the first two true future frames. The digits do not overlap in any
displayed true frame.}
\label{fig:q2-f5}
\end{figure}

%Q29
\qa{9} Model A falls apart at 2$\times$ already. Its $\Ex2$ jumps from 0.0189 at 1$\times$ to 0.0439
($2.3\times$), slightly worse than predicting nothing (Blank, 0.0426), and rises to 0.0513 at 3$\times$.
It stays below Copy-last (0.0703) only because a blurred guess is penalised less than a sharp misplaced digit.
This does \emph{not} agree with my reach estimate. At 2$\times$ the median displacement (6.6\,px) is
within the 8\,px reach, so the architecture could carry the information far enough, yet the model fails
there. The estimate measures what the architecture \emph{could} do with some weights. It ignores what the
trained weights \emph{do}, and they were trained only on displacements of at most about 4.7\,px. Within one step,
the filters have learned to move a digit by about one cell or less, and to infer velocity only from small
offsets between the new frame and the remembered position. Nothing in training asked the second hop to
carry a larger shift. The reach is an upper bound, and here the learned behaviour falls well short of it.

Measuring blob centroids in Figure~\ref{fig:q2-f5} (\texttt{q2/positions.py}) makes the failure concrete.
At 1$\times$ both predicted digits sit within 1\,px of the truth at $h=1$; at $h=2$ the 1 misses its bounce
off the left wall and keeps going left. At 2$\times$ the digits \emph{lag}: they move in the right
direction but only 1.8--2.5\,px per step, against 7.3--7.6\,px in the truth, so they fall further behind at
every horizon. At 3$\times$ they \emph{stall}: both predicted digits stay within about 1\,px of their last
input position, while the truth moves 5.8 and 11.2\,px. The prediction is essentially a blurred Copy-last.
At both speeds the digits are also \emph{blurred}: the blobs above intensity 0.2 cover 1.6--1.9 times
as many pixels as the true digits, at a lower peak intensity (about 0.75). Neither digit is duplicated or vanishes within $h\le2$.
The lag suggests that the model applies the speed it learned in training instead of measuring the speed
in front of it. The stall at 3$\times$ suggests that it can no longer match a digit to its previous position
(the shift is beyond what it learned, and beyond its reach), so it falls back to persistence. The blur is
the MSE-optimal hedge when the position is uncertain: spread intensity over the possible locations.

\qa{10} No, not from these results alone. At 3$\times$ both explanations hold at once: the motion
(about 10\,px/frame) was never seen in training \emph{and} it exceeds the 8\,px reach, so the failure
cannot be attributed to either one. 2$\times$ suggests that unfamiliarity already matters, because the
model fails there even though the motion is within its reach. But it does not show that reach is
irrelevant at 3$\times$. Two experiments, each changing one factor, separate them.
(i) \emph{Change only the data}: retrain the same architecture with speed augmentation (model C). If
unfamiliar motion is the cause, C should recover at both 2$\times$ and 3$\times$; if reach is the cause,
C should recover at 2$\times$ (within reach) but stay poor at 3$\times$, because no amount of data lets
information travel further than the architecture allows within one step.
(ii) \emph{Change only the reach}: train an architecture with a larger reach on exactly the same
speed-augmented data (model D). If D fixes 3$\times$ where C could not, the remaining failure was
architectural reach.

%======================================================================
%======================================================================
\section{What Can Data Fix?}
%======================================================================
Models B, C, D and the optional D+B are trained with exactly the protocol of model A
(\path{q2/train.py}, configurations \texttt{B}, \texttt{C}, \texttt{D}, \texttt{DB}). B draws
$g\in G$ uniformly for every training window and applies it to input and target. C, D and D+B draw
$k\in\{1,2,3\}$ uniformly and then the start frame uniformly from the allowed starts for that $k$; D+B
also draws $g$. Training videos with only two existing targets contribute equally, through the masked loss.

\qa{11} \emph{Included exactly as first written. It was committed (\texttt{Question2/predictions/Q2\_11.md},
commit \texttt{20f3c69}, 2026-10-29 04:28 PKT) after the Section~3 results and before any of B--D was trained
(B seed 0 started at 04:29 PKT).}

\begin{quote}\small
Model B (symmetry augmentation) will shrink C\_M and C\_R several-fold and bring E\_M and
E\_R down to its E(5) at 1x, but will not change the speed failure. Model C (speed augmentation) will
largely fix 2x, where the motion (about 6.6 px per step) is still within model A's 8 px reach, but only
partly fix 3x (about 10 px per step, beyond that reach), and will not change the consistency gaps.
Model D (reach 16 px plus speed augmentation) will also fix 3x, bringing E(2) at 3x close to its value at
2x at a small cost in E(5) at 1x from the coarser 8x8 grid, while its consistency gaps stay nonzero like
those of A and C.
\end{quote}

\begin{table}[H]
\centering
\caption{Models A--D and the optional D+B on the test set. $C_T$ and $E_T$ at 1$\times$; $\Ex2$ uses the
first two horizons. Training time is the mean wall-clock time per epoch over 30 epochs; inference time is
for all 1,000 test videos at batch 32 (hardware in the introduction to this question).}
\label{tab:q2-t5}
\resizebox{\linewidth}{!}{\begin{tabular}{lccccc}
\toprule
 & A & B & C & D & D+B \\
\midrule
$E^{(5)}$ 1$\times$ & 0.0224 (0.0228, 0.0219) & 0.0224 (0.0232, 0.0215) & 0.0271 (0.0276, 0.0266) & 0.0258 (0.0260, 0.0257) & 0.0266 (0.0265, 0.0267) \\
$C_M$ & 0.00156 (0.00148, 0.00164) & 0.00133 (0.00140, 0.00126) & $9.85\times10^{-4}$ ($7.42\times10^{-4}$, 0.00123) & 0.00117 (0.00104, 0.00131) & 0.00137 (0.00168, 0.00106) \\
$E_M$ & 0.0230 (0.0234, 0.0226) & 0.0224 (0.0232, 0.0216) & 0.0273 (0.0278, 0.0269) & 0.0261 (0.0262, 0.0260) & 0.0266 (0.0266, 0.0267) \\
$C_R$ & 0.00275 (0.00243, 0.00307) & 0.00175 (0.00187, 0.00164) & 0.00230 (0.00161, 0.00300) & 0.00253 (0.00239, 0.00268) & 0.00145 (0.00162, 0.00128) \\
$E_R$ & 0.0240 (0.0243, 0.0238) & 0.0224 (0.0231, 0.0217) & 0.0283 (0.0286, 0.0280) & 0.0273 (0.0274, 0.0272) & 0.0266 (0.0265, 0.0267) \\
$E^{(2)}$ 1$\times$ & 0.0189 (0.0193, 0.0185) & 0.0186 (0.0191, 0.0181) & 0.0233 (0.0236, 0.0229) & 0.0236 (0.0237, 0.0236) & 0.0243 (0.0240, 0.0245) \\
$E^{(2)}$ 2$\times$ & 0.0439 (0.0428, 0.0450) & 0.0453 (0.0448, 0.0458) & 0.0291 (0.0294, 0.0288) & 0.0260 (0.0260, 0.0260) & 0.0265 (0.0265, 0.0265) \\
$E^{(2)}$ 3$\times$ & 0.0513 (0.0505, 0.0522) & 0.0525 (0.0516, 0.0534) & 0.0336 (0.0338, 0.0335) & 0.0292 (0.0292, 0.0291) & 0.0295 (0.0295, 0.0296) \\
Trainable parameters & 480,784 & 480,784 & 480,784 & 594,496 & 594,496 \\
Training time (s/epoch) & 207 (210, 203) & 204 (204, 205) & 209 (216, 202) & 59.7 (60.0, 59.4) & 59.7 (59.8, 59.5) \\
Inference time (s, 1000 videos) & 7.83 (7.83, 7.83) & 7.92 (7.96, 7.88) & 7.96 (7.96, 7.96) & 2.21 (2.20, 2.21) & 2.19 (2.18, 2.19) \\
\bottomrule
\end{tabular}
}
\end{table}

\qa{12} \emph{Model D} is model A with a coarser fold: \texttt{pixel\_unshuffle}$(x,8)$ turns a
$1\times64\times64$ frame into $64\times8\times8$. Everything else is unchanged: two peephole-free
ConvLSTM layers with 64 hidden channels and $3\times3$ kernels (stride 1, padding 1), a $1\times1$
convolution from 64 to 64 channels, \texttt{pixel\_shuffle}$(\cdot,8)$, a sigmoid, the same 9-step
closed-loop rollout without teacher forcing, and the same training protocol (with speed augmentation, as
in its row of the model table). Its parameters are $2\times\big[(64+64)\cdot4\cdot64\cdot9+256\big]$ for the
two ConvLSTM layers plus $64\cdot64+64$ for the output: \textbf{594,496}, i.e.\ $1.24\times$ model A and
under the 721,176 limit.
\emph{Reach.} The same path argument as in Q2.3 applies, but one cell of the $8\times8$ grid is now
8 original pixels. From $u_\tau$, or from $h^{(1)}_{\tau-1}$, a path crosses layer~1's and layer~2's
$3\times3$ convolutions: 2 cells $=16$\,px. From $h^{(2)}_{\tau-1}$ it crosses only layer~2's: 8\,px.
The $1\times1$ convolution and the (un)shuffles move nothing. The reach is therefore \textbf{16\,px per
step} along a row or column ($16\sqrt2\approx22.6$\,px diagonally), strictly larger than model A's 8\,px.
It covers the 3$\times$ median of 10\,px/frame and its 95th percentile of 13\,px.
\emph{Why this design.} A bigger patch doubles the reach of every convolution without adding layers, and
the cost per step goes down (an $8\times8$ grid instead of $16\times16$). A dilation of 2 would give the
same 16\,px at A's parameter count, but each hop could then reach only even cell offsets, leaving gaps in
the neighbourhood. $5\times5$ kernels at 64 channels would exceed the parameter limit. The price of the
coarser fold is spatial detail: each state cell must now represent and decode an $8\times8$ patch from 64
channels, which may cost sharpness at 1$\times$. The design was fixed a priori from this reach calculation.
No candidate designs were compared, so neither the validation nor the test split influenced it.
\emph{Initialisation.} With PyTorch's default initialisation (the one model A uses), seed~0 of model D
never left a constant output: from the first epoch it predicted the same faint grey value (about 0.05) at every
pixel, whatever the input, and its validation loss stayed at 0.03947 for all 30 epochs. Seed~1
trained normally (validation loss 0.0267 after 30 epochs). I therefore changed model D's initialisation,
for both seeds and for D+B: Xavier-uniform weights in every convolution, zero biases, and forget-gate
biases of 1 in both ConvLSTM cells, so that cell state and gradients pass through the recurrence from
the first step (\texttt{xavier\_forget1\_} in \texttt{q2/models.py}). All four D and D+B runs were then
retrained from scratch with seeds 0 and 1; every D and D+B number in this report comes from them. The decision used
only the training curve of the failed run, not the test split. The training logs of the first runs are in \path{Question2/logs/old_init/}. With the $8\times8$ grid, an epoch also takes
about 60\,s instead of about 205\,s (Table~\ref{tab:q2-t5}): the convolutions run on 4 times fewer
positions, with only slightly more input channels per position in layer~1.

\qa{13} The three failures from Testing the Baseline get three different labels.

\emph{The consistency gaps under $M$ and $R$: appearance.} Showing the model transformed digits fixes them,
and showing it more speeds does not. Model B (symmetry augmentation, nothing else changed) brings both
transformed errors down to its ordinary error: $E_M=E_R=\Ex5=0.0224$ at 1$\times$, against
$E_M=0.0230$ and $E_R=0.0240$ for A. Its $C_R$ drops by a third ($0.00275\to0.00175$) and $C_M$ by a sixth
($0.00156\to0.00133$). Speed augmentation does not do this: for C and D, $E_R$ stays 0.0012--0.0015 above $\Ex5$,
as for A (0.0016), and $C_R$ stays at 0.0023--0.0025. Adding symmetry augmentation to D (D+B) closes the gap
again: $E_R=E_M=\Ex5=0.0266$ and $C_R=0.00145$. By Q2.1--Q2.2, every mirrored or rotated motion already
occurs in training, so what B adds is mainly how transformed digits look: the label is appearance.

\emph{The failure at 2$\times$: unfamiliar motion.} Model C has model A's architecture and reach and
differs only in its training speeds. It lowers $\Ex2$ at 2$\times$ from 0.0439 to 0.0291, much closer to
its own 0.0233 at 1$\times$ than to A's failure. The 2$\times$ motion is within A's 8\,px reach, so
the only thing that changed is that the model has now seen it. In the Figure~5 video, model C moves the
digits 4.6\,px at $h=1$ instead of model A's 1.8--2.5\,px (true 7.3--7.6\,px).

\emph{The failure at 3$\times$: mostly unfamiliar motion, partly architectural reach.} Model C also lowers
$\Ex2$ at 3$\times$, from 0.0513 to 0.0336. So most of the 3$\times$ failure was unfamiliar motion too:
model C does not ``still do badly'' at 3$\times$. Model D changes only the reach relative to C and lowers
it further to 0.0292, and it also lowers 2$\times$ to 0.0260. Relative to its own 1$\times$ error, model D
loses 24\,\% at 3$\times$ ($0.0236\to0.0292$) where C loses 44\,\% ($0.0233\to0.0336$). The part of the
3$\times$ failure that remains in C is therefore architectural reach. More data cannot carry information
further: the reach is a property of the computation graph. A predicted pixel at step $\tau+1$ is a
function of only the input pixels within 8\,px (for A and C), whatever the weights are, and training data
can only change the weights. A digit that moves 10\,px per frame arrives where nothing in the previous frame
could tell the model it was coming. C can only partly work around this: a digit is about 20\,px wide, so
most of its new position still lies within 8\,px of its old pixels, and the recurrent state already
holds its velocity from earlier frames. C can therefore smear the digit in the direction of motion, but it
cannot place the leading edge exactly. The positions in the Figure~5 video show both effects. At 2$\times$,
$h=2$, model D places the two digits 0.5 and 1.3\,px from the truth, while C is 3.5 and 7.7\,px off and
its blobs fade (45 and 61 pixels above 0.2, against 83--86 in the truth). At 3$\times$, $h=1$, both C and D move the lower digit 8.0--8.4\,px (true 11.2\,px), and both miss the
upper digit's bounce off the top wall; by $h=2$ C merges the two digits into one blob by $h=2$, whereas D keeps them apart.
One caveat: D also coarsens the grid, so the reach is not the only thing that differs between C and D.
D is also better than C at 1$\times$ (0.0258 vs 0.0271), where reach does not limit either model.

\qa{14} Yes. Speed augmentation costs accuracy on the original 1$\times$ clips, and symmetry
augmentation does not. $\Ex5$ at 1$\times$ is 0.0224 for A and 0.0224 for B, but 0.0271 for C
(+21\,\%), 0.0258 for D (+15\,\%) and 0.0266 for D+B (+19\,\%), consistently in both seeds; $\Ex2$ at
1$\times$ rises from 0.0189 to 0.0233--0.0243. The cost comes from the speed augmentation, not from the
coarser grid of D: D is better than C at 1$\times$. Only a third of the training windows are now at 1$\times$,
and predictions at 1$\times$ are visibly blurrier (compare the top rows of Figure~\ref{fig:q2-f5}, model A, and
Figure~\ref{fig:q2-f6}b, the final model). A model that must cover
three speeds also has more ways to hedge between them. Averaged over the three speeds, $\Ex2$ falls from
0.0380 (A) to 0.0263 (D), a 31\,\% improvement, and at 2$\times$ and 3$\times$ model A is no better than
predicting a blank frame. If the forecaster only ever sees clips at the original speed, the change is not
worth it: A or B is the better model. As a forecaster of motion in general, which is the point of this
question, it is worth it: D trades a 15\,\% loss at 1$\times$ for a 41--43\,\% gain at 2$\times$ and
3$\times$.

\emph{How the Q2.11 prediction fared.} Mostly right, with two misses. B did bring $E_M$ and $E_R$ down to its
$\Ex5$ at 1$\times$ and did not change the speed failure, but it shrank $C_M$ and $C_R$ by 15\,\% and
36\,\%, not ``several-fold''. C largely fixed 2$\times$ and partly fixed 3$\times$, as predicted; its
consistency gaps stayed of the same size as A's ($C_M$ even fell somewhat). D brought $\Ex2$ at 3$\times$
(0.0292) close to 2$\times$ (0.0260), with gaps still nonzero, as predicted. But its cost at 1$\times$ came
from the speed augmentation it shares with C, not from the coarser grid.

%======================================================================
\section{The Final Model}
%======================================================================
\emph{The symmetrised forecaster is exactly consistent.} Let $T\in G$, and recall that every $g\in G$
acts on frames as a pixel permutation, so it is linear and commutes with averaging. Then
\[
\bar f(TX)=\frac18\sum_{g\in G} g^{-1} f(gTX)
=\frac18\sum_{g'\in G}T\,g'^{-1} f(g'X)=T\Big(\frac18\sum_{g'\in G}g'^{-1}f(g'X)\Big)=T\bar f(X),
\]
substituting $g'=gT$. As $g$ runs over $G$, so does $g'$, because $G$ is a group; and
$g^{-1}=(g'T^{-1})^{-1}=Tg'^{-1}$. Hence $\bar f(T(X))=T(\bar f(X))$ for every input, and
$C_T=\mathrm{MSE}\big(\bar f(T(X)),T(\bar f(X))\big)=0$ for $T=M$ and $T=R$ (and every other element of
$G$), whatever $f$ is. In floating point, the two sides add the same eight terms in a different order,
so the measured gaps are rounding errors only (Table~\ref{tab:q2-t6}). The implementation is
\texttt{Symmetrised} in \texttt{q2/models.py}; inference timing includes the eight transformations, the
eight forward passes, the inverse transformations and the averaging.

\emph{Choice of final model: symmetrised D.} The default is model D, and Table~\ref{tab:q2-t5} does not
make a case for D+B. Symmetry augmentation on top of D buys learned consistency ($C_R$ $0.00253\to0.00145$,
$E_R$ $0.0273\to0.0266$ at 1$\times$), but it costs ordinary accuracy at 1$\times$ ($\Ex5$
$0.0258\to0.0266$, in both seeds) and gains nothing at 2$\times$ or 3$\times$ (0.0260 vs 0.0265,
0.0292 vs 0.0295). Once wrapped, any model is exactly consistent, so what D+B adds is redundant.
Table~\ref{tab:q2-t6} confirms this: $\bar f$(D) and $\bar f$(D+B) differ by at most 0.0002 in every
error column, with $\bar f$(D) slightly ahead at 1$\times$ (0.0258 vs 0.0260), at the same parameter count and cost.

\begin{table}[H]
\centering
\caption{Symmetrised forecasters $\bar f$ built from the seed-0 and seed-1 models B, D and D+B. Training time is
that of the underlying learned model (symmetrisation needs no training); inference time includes all
eight transformed forward passes.}
\label{tab:q2-t6}
\resizebox{\linewidth}{!}{\begin{tabular}{lccc}
\toprule
 & $\bar f$(B) & $\bar f$(D) & $\bar f$(D+B) \\
\midrule
$E^{(5)}$ 1$\times$ & 0.0217 (0.0224, 0.0209) & 0.0258 (0.0260, 0.0256) & 0.0260 (0.0258, 0.0262) \\
$C_M$ & $9.37\times10^{-17}$ ($9.44\times10^{-17}$, $9.31\times10^{-17}$) & $6.41\times10^{-17}$ ($5.82\times10^{-17}$, $7.00\times10^{-17}$) & $6.62\times10^{-17}$ ($7.21\times10^{-17}$, $6.02\times10^{-17}$) \\
$E_M$ & 0.0217 (0.0224, 0.0209) & 0.0258 (0.0260, 0.0256) & 0.0260 (0.0258, 0.0262) \\
$C_R$ & $7.55\times10^{-17}$ ($7.61\times10^{-17}$, $7.49\times10^{-17}$) & $5.13\times10^{-17}$ ($4.66\times10^{-17}$, $5.61\times10^{-17}$) & $5.30\times10^{-17}$ ($5.82\times10^{-17}$, $4.79\times10^{-17}$) \\
$E_R$ & 0.0217 (0.0224, 0.0209) & 0.0258 (0.0260, 0.0256) & 0.0260 (0.0258, 0.0262) \\
$E^{(2)}$ 1$\times$ & 0.0181 (0.0185, 0.0177) & 0.0237 (0.0238, 0.0236) & 0.0239 (0.0237, 0.0241) \\
$E^{(2)}$ 2$\times$ & 0.0444 (0.0439, 0.0449) & 0.0260 (0.0260, 0.0260) & 0.0261 (0.0261, 0.0261) \\
$E^{(2)}$ 3$\times$ & 0.0512 (0.0504, 0.0520) & 0.0289 (0.0291, 0.0288) & 0.0290 (0.0290, 0.0291) \\
Trainable parameters & 480,784 & 594,496 & 594,496 \\
Training time (s/epoch) & 204 (204, 205) & 59.7 (60.0, 59.4) & 59.7 (59.8, 59.5) \\
Inference time (s, 1000 videos) & 64.3 (64.4, 64.2) & 17.7 (17.8, 17.6) & 17.8 (17.7, 17.8) \\
\bottomrule
\end{tabular}
}
\end{table}

\begin{figure}[H]
\centering
\includegraphics[width=0.92\linewidth]{fig6a_final_M_9269.png}\\[2pt]
\includegraphics[width=0.92\linewidth]{fig6a_final_M_9256.png}
\caption[Final model, Figure 3 and Figure 5 repeated]{(a) Figure~\ref{fig:q2-f3} repeated with the final
model (symmetrised D, seed 0) on the same videos: $T=M$, videos 9269 (top) and 9256 (bottom); continued
below. Rows as in Figure~\ref{fig:q2-f3}. GIFs: \path{Question2/gifs/fig6a_final_M_video9269.gif},
\path{Question2/gifs/fig6a_final_M_video9256.gif}, \path{Question2/gifs/fig6a_final_R_video9983.gif},
\path{Question2/gifs/fig6a_final_R_video9733.gif}.}
\label{fig:q2-f6}
\end{figure}
\begin{figure}[H]
\ContinuedFloat
\centering
\includegraphics[width=0.92\linewidth]{fig6a_final_R_9983.png}\\[2pt]
\includegraphics[width=0.92\linewidth]{fig6a_final_R_9733.png}
\caption{(continued) (a) $T=R$, videos 9983 (top) and 9733 (bottom). Rows 3 and 4 now coincide, because
$\bar f(R(X))=R(\bar f(X))$.}
\end{figure}
\begin{figure}[H]
\ContinuedFloat
\centering
\includegraphics[width=\linewidth]{fig6b_speed_final.png}
\caption{(continued) (b) Figure~\ref{fig:q2-f5} repeated with the final model (symmetrised D, seed 0) on test
video 9021 at 1$\times$, 2$\times$ and 3$\times$ ($t=1$).}
\end{figure}

Measured in the same way as for model A in Q2.9 (\texttt{python -m q2.positions sym-D}), the final model
follows the digits of Figure~\ref{fig:q2-f6}b. At 1$\times$ both digits are within 1.3\,px of the truth at
$h=1$ (model A: within 0.5\,px). At $h=2$ it also misses the 1's bounce off the left wall, but by less
(4.1\,px, model A 6.2\,px). At 2$\times$ it is within 1.5\,px at $h=1$ and within 1.8\,px at $h=2$,
where model A was 4.9--8.8\,px off. At 3$\times$, $h=1$, it moves the lower digit 7.8\,px of the true
11.2\,px (4.7\,px off) and the upper digit, which bounces off the top wall, ends 3.2\,px off. Its blobs
are larger than A's at 1$\times$, but smaller at 2$\times$ and 3$\times$.

\qa{15} \emph{Consistency.} Built-in symmetry is exact and learned symmetry is not. Symmetrised B has
$C_M=9.4\times10^{-17}$ and $C_R=7.6\times10^{-17}$, rounding error in single precision, against
$C_M=0.00133$ and $C_R=0.00175$ for model B itself, which is only somewhat lower than model A. The
argument above shows why: the group average makes $\bar f(T(X))=T(\bar f(X))$ an identity that holds for
every input and every set of weights. Augmentation only rewards the network for being \emph{right} on
transformed clips, on average over the training data. Nothing ties its prediction for $T(X)$ to its
prediction for $X$, and with finite capacity and training it reaches neither on unseen videos.
\emph{Accuracy.} On transformed clips, $E_M=E_R=\Ex5$ for every wrapped model, by construction. On the
original clips, wrapping did not cost accuracy anywhere: $\Ex5$ at 1$\times$ goes from 0.0224 to 0.0217
for B ($-3\,\%$, both seeds), stays at 0.0258 for D, and goes from 0.0266 to 0.0260 for D+B. For D, the
transformed errors fall from 0.0261 ($M$) and 0.0273 ($R$) to 0.0258. $\Ex2$ at 2$\times$ and 3$\times$
changes by at most 0.0009.
\emph{Inference cost.} Wrapping costs eight forward passes: 7.92\,s to 64.3\,s for B and 2.21\,s to 17.7\,s
for D on the 1,000 test videos (about $8\times$), with no training.
\emph{Blur.} The average of eight predictions is blurred wherever they disagree. Because MSE is convex,
the wrapped forecaster's error is never larger than the mean of the eight transformed errors $E_g$ of
$f$. It can, however, be larger than $f$'s own error on the original clips, if $f$ is much better on
untransformed clips than on transformed ones, i.e.\ when $E_T\gg\Ex5$. That is the case for a model trained without
symmetry augmentation whose motion rule depends on appearance (model A, $E_R-\Ex5=0.0016$), and for
videos where the eight views give different positions. Here it did not hurt: the gaps $E_T-\Ex5$ of B
and D are small, and $\Ex5$ at 1$\times$ improved or stayed the same, so the blurring averages out
errors more than it smears correct predictions. The largest change for the worse in any column is $+0.0001$ ($\Ex2$ at
1$\times$ for D).
\qa{16} Only partly. In Q2.7 the failure on video 9983 ($T=R$) had two parts: the two accounts of
the motion disagreed (rows 3 and 4 of Figure~\ref{fig:q2-f3}c), and neither one separated the overlapping
3 and 0. The first part is fixed. In Figure~\ref{fig:q2-f6}a (video 9983), rows 3 and 4 are identical, and
the video's $C_R$ is 0 up to rounding (model A, seed 0: 0.00502). The change that fixed it is the symmetrisation; symmetry
augmentation alone (model B) only reduced it to 0.00216. The second part is not fixed. The final model still
predicts one tall blob that drifts slightly downwards and fades, while the true digits separate and move
to the lower left. On this video it is no more accurate than model A: its error is 0.0338 on both the
original and the rotated clip, against 0.0292 and 0.0335 for model A. The speed augmentation made this
1$\times$ video worse: model C gives 0.0361 and model D 0.0330 on the original clip. This matches Q2.14.
None of the changes addresses what makes this video hard: two overlapping, curly digits whose appearance
does not tell the model where each one is.
\qa{17} Partly. The final model learned a motion rule that transfers across speeds and, by construction,
across all eight symmetries. It is not a persistence or blank predictor: at 3$\times$ its $\Ex2$ is 0.0289,
against 0.0703 for Copy-last and 0.0426 for Blank, and at 2$\times$ its digits move almost the true
distance (Figure~\ref{fig:q2-f6}b). To convince a skeptic, I would measure centroid tracking errors over
the whole test set, per digit, horizon and speed, including wall bounces, against Copy-last and
constant-velocity extrapolation. I would also test an unseen speed such as 4$\times$, where a learned rule
should degrade gracefully instead of stalling as model A did at 3$\times$. One test it has not passed is
keeping digits intact: overlapping digits merge into one blob (video 9983), bounces are missed at
3$\times$, and by $h=5$ every prediction is a blurred shape rather than a digit.

